\chapter{Quyền im lặng}
\markboth{Quyền im lặng}{The Right to Silence}


\section{Em không bắt buộc phải trả lời ngay; em có quyền nghĩ trước.}
\begin{truyen}{Em không bắt buộc phải trả lời ngay; em có quyền nghĩ trước.}{Em không bắt buộc phải trả lời ngay; em có quyền n}
\chuhoa{C}{âu thầy nói câu ấy. Cả lớp im lặng — không ai nói, mỗi người đang nghĩ về điều vừa nghe.}
\textit{(The teacher said that. The class fell silent — no one spoke; everyone was thinking about what they had just heard.)}
\end{truyen}

\begin{baihoc}
Im lặng sau câu nói đúng lúc là dấu hiệu của tư duy.
\textit{(Silence after the right sentence is a sign of thinking.)}
\end{baihoc}

\begin{ghinhoanh}
\textbf{Short English to remember:}

Some sentences make the whole class think in silence.

\end{ghinhoanh}

\begin{tuhocnhanh}
1. Câu này khiến em nghĩ đến điều gì? \textit{What does this make you think of?}

2. Em có thể dùng câu này trong tình huống nào? \textit{When could you use this sentence?}
\end{tuhocnhanh}
\ngancach


\section{Im lặng là khoảng trống để ý tưởng lớn lên.}
\begin{truyen}{Im lặng là khoảng trống để ý tưởng lớn lên.}{Im lặng là khoảng trống để ý tưởng lớn lên.}
\chuhoa{C}{âu thầy nói câu ấy. Cả lớp im lặng — không ai nói, mỗi người đang nghĩ về điều vừa nghe.}
\textit{(The teacher said that. The class fell silent — no one spoke; everyone was thinking about what they had just heard.)}
\end{truyen}

\begin{baihoc}
Im lặng sau câu nói đúng lúc là dấu hiệu của tư duy.
\textit{(Silence after the right sentence is a sign of thinking.)}
\end{baihoc}

\begin{ghinhoanh}
\textbf{Short English to remember:}

Some sentences make the whole class think in silence.

\end{ghinhoanh}

\begin{tuhocnhanh}
1. Câu này khiến em nghĩ đến điều gì? \textit{What does this make you think of?}

2. Em có thể dùng câu này trong tình huống nào? \textit{When could you use this sentence?}
\end{tuhocnhanh}
\ngancach


\section{Câu hỏi hay thường được theo sau bởi im lặng — vì mọi người đang nghĩ.}
\begin{truyen}{Câu hỏi hay thường được theo sau bởi im lặng — vì mọi người ...}{Câu hỏi hay thường được theo sau bởi im lặng — vì }
\chuhoa{C}{âu thầy nói câu ấy. Cả lớp im lặng — không ai nói, mỗi người đang nghĩ về điều vừa nghe.}
\textit{(The teacher said that. The class fell silent — no one spoke; everyone was thinking about what they had just heard.)}
\end{truyen}

\begin{baihoc}
Im lặng sau câu nói đúng lúc là dấu hiệu của tư duy.
\textit{(Silence after the right sentence is a sign of thinking.)}
\end{baihoc}

\begin{ghinhoanh}
\textbf{Short English to remember:}

Some sentences make the whole class think in silence.

\end{ghinhoanh}

\begin{tuhocnhanh}
1. Câu này khiến em nghĩ đến điều gì? \textit{What does this make you think of?}

2. Em có thể dùng câu này trong tình huống nào? \textit{When could you use this sentence?}
\end{tuhocnhanh}
\ngancach


\section{Người biết im lặng đúng lúc thường là người nói có sức nặng.}
\begin{truyen}{Người biết im lặng đúng lúc thường là người nói có sức nặng.}{Người biết im lặng đúng lúc thường là người nói có}
\chuhoa{C}{âu thầy nói câu ấy. Cả lớp im lặng — không ai nói, mỗi người đang nghĩ về điều vừa nghe.}
\textit{(The teacher said that. The class fell silent — no one spoke; everyone was thinking about what they had just heard.)}
\end{truyen}

\begin{baihoc}
Im lặng sau câu nói đúng lúc là dấu hiệu của tư duy.
\textit{(Silence after the right sentence is a sign of thinking.)}
\end{baihoc}

\begin{ghinhoanh}
\textbf{Short English to remember:}

Some sentences make the whole class think in silence.

\end{ghinhoanh}

\begin{tuhocnhanh}
1. Câu này khiến em nghĩ đến điều gì? \textit{What does this make you think of?}

2. Em có thể dùng câu này trong tình huống nào? \textit{When could you use this sentence?}
\end{tuhocnhanh}
\ngancach


\section{Học không chỉ là nghe và trả lời — học còn là im lặng và suy nghĩ.}
\begin{truyen}{Học không chỉ là nghe và trả lời — học còn là im lặng và suy...}{Học không chỉ là nghe và trả lời — học còn là im l}
\chuhoa{C}{âu thầy nói câu ấy. Cả lớp im lặng — không ai nói, mỗi người đang nghĩ về điều vừa nghe.}
\textit{(The teacher said that. The class fell silent — no one spoke; everyone was thinking about what they had just heard.)}
\end{truyen}

\begin{baihoc}
Im lặng sau câu nói đúng lúc là dấu hiệu của tư duy.
\textit{(Silence after the right sentence is a sign of thinking.)}
\end{baihoc}

\begin{ghinhoanh}
\textbf{Short English to remember:}

Some sentences make the whole class think in silence.

\end{ghinhoanh}

\begin{tuhocnhanh}
1. Câu này khiến em nghĩ đến điều gì? \textit{What does this make you think of?}

2. Em có thể dùng câu này trong tình huống nào? \textit{When could you use this sentence?}
\end{tuhocnhanh}
\ngancach


\section{Im lặng không phải từ chối tham gia — mà là tham gia bằng cách nghĩ.}
\begin{truyen}{Im lặng không phải từ chối tham gia — mà là tham gia bằng cá...}{Im lặng không phải từ chối tham gia — mà là tham g}
\chuhoa{C}{âu thầy nói câu ấy. Cả lớp im lặng — không ai nói, mỗi người đang nghĩ về điều vừa nghe.}
\textit{(The teacher said that. The class fell silent — no one spoke; everyone was thinking about what they had just heard.)}
\end{truyen}

\begin{baihoc}
Im lặng sau câu nói đúng lúc là dấu hiệu của tư duy.
\textit{(Silence after the right sentence is a sign of thinking.)}
\end{baihoc}

\begin{ghinhoanh}
\textbf{Short English to remember:}

Some sentences make the whole class think in silence.

\end{ghinhoanh}

\begin{tuhocnhanh}
1. Câu này khiến em nghĩ đến điều gì? \textit{What does this make you think of?}

2. Em có thể dùng câu này trong tình huống nào? \textit{When could you use this sentence?}
\end{tuhocnhanh}
\ngancach


\section{Thầy chấp nhận im lặng — vì thầy biết suy nghĩ cần thời gian.}
\begin{truyen}{Thầy chấp nhận im lặng — vì thầy biết suy nghĩ cần thời gian...}{Thầy chấp nhận im lặng — vì thầy biết suy nghĩ cần}
\chuhoa{C}{âu thầy nói câu ấy. Cả lớp im lặng — không ai nói, mỗi người đang nghĩ về điều vừa nghe.}
\textit{(The teacher said that. The class fell silent — no one spoke; everyone was thinking about what they had just heard.)}
\end{truyen}

\begin{baihoc}
Im lặng sau câu nói đúng lúc là dấu hiệu của tư duy.
\textit{(Silence after the right sentence is a sign of thinking.)}
\end{baihoc}

\begin{ghinhoanh}
\textbf{Short English to remember:}

Some sentences make the whole class think in silence.

\end{ghinhoanh}

\begin{tuhocnhanh}
1. Câu này khiến em nghĩ đến điều gì? \textit{What does this make you think of?}

2. Em có thể dùng câu này trong tình huống nào? \textit{When could you use this sentence?}
\end{tuhocnhanh}
\ngancach


\section{Điều em nghĩ trong im lặng có thể quan trọng hơn điều em nói trong ồn ào.}
\begin{truyen}{Điều em nghĩ trong im lặng có thể quan trọng hơn điều em nói...}{Điều em nghĩ trong im lặng có thể quan trọng hơn đ}
\chuhoa{C}{âu thầy nói câu ấy. Cả lớp im lặng — không ai nói, mỗi người đang nghĩ về điều vừa nghe.}
\textit{(The teacher said that. The class fell silent — no one spoke; everyone was thinking about what they had just heard.)}
\end{truyen}

\begin{baihoc}
Im lặng sau câu nói đúng lúc là dấu hiệu của tư duy.
\textit{(Silence after the right sentence is a sign of thinking.)}
\end{baihoc}

\begin{ghinhoanh}
\textbf{Short English to remember:}

Some sentences make the whole class think in silence.

\end{ghinhoanh}

\begin{tuhocnhanh}
1. Câu này khiến em nghĩ đến điều gì? \textit{What does this make you think of?}

2. Em có thể dùng câu này trong tình huống nào? \textit{When could you use this sentence?}
\end{tuhocnhanh}
\ngancach


\section{Lớp học tốt là nơi im lặng được tôn trọng.}
\begin{truyen}{Lớp học tốt là nơi im lặng được tôn trọng.}{Lớp học tốt là nơi im lặng được tôn trọng.}
\chuhoa{C}{âu thầy nói câu ấy. Cả lớp im lặng — không ai nói, mỗi người đang nghĩ về điều vừa nghe.}
\textit{(The teacher said that. The class fell silent — no one spoke; everyone was thinking about what they had just heard.)}
\end{truyen}

\begin{baihoc}
Im lặng sau câu nói đúng lúc là dấu hiệu của tư duy.
\textit{(Silence after the right sentence is a sign of thinking.)}
\end{baihoc}

\begin{ghinhoanh}
\textbf{Short English to remember:}

Some sentences make the whole class think in silence.

\end{ghinhoanh}

\begin{tuhocnhanh}
1. Câu này khiến em nghĩ đến điều gì? \textit{What does this make you think of?}

2. Em có thể dùng câu này trong tình huống nào? \textit{When could you use this sentence?}
\end{tuhocnhanh}
\ngancach


\section{Suy nghĩ sâu thường bắt đầu từ một khoảnh khắc im lặng.}
\begin{truyen}{Suy nghĩ sâu thường bắt đầu từ một khoảnh khắc im lặng.}{Suy nghĩ sâu thường bắt đầu từ một khoảnh khắc im }
\chuhoa{C}{âu thầy nói câu ấy. Cả lớp im lặng — không ai nói, mỗi người đang nghĩ về điều vừa nghe.}
\textit{(The teacher said that. The class fell silent — no one spoke; everyone was thinking about what they had just heard.)}
\end{truyen}

\begin{baihoc}
Im lặng sau câu nói đúng lúc là dấu hiệu của tư duy.
\textit{(Silence after the right sentence is a sign of thinking.)}
\end{baihoc}

\begin{ghinhoanh}
\textbf{Short English to remember:}

Some sentences make the whole class think in silence.

\end{ghinhoanh}

\begin{tuhocnhanh}
1. Câu này khiến em nghĩ đến điều gì? \textit{What does this make you think of?}

2. Em có thể dùng câu này trong tình huống nào? \textit{When could you use this sentence?}
\end{tuhocnhanh}
\ngancach
